\documentclass[11pt]{article}

\usepackage[margin=0.8in]{geometry}
\usepackage{amsmath,amssymb}
\usepackage{enumitem}

\setlength{\parindent}{0pt}
\setlength{\parskip}{5pt}

\begin{document}

{\large \textbf{LLMs for Program Repair}}

\vspace{7pt}

Normal APR tries to learn transformations like

\[
\text{buggy code} \rightarrow \text{fixed code}
\]

But with LLMs, another idea is there:

give the model the surrounding code and directly ask
it to predict the correct code.

So model uses context of the buggy program.

\vspace{7pt}

\textbf{code infilling}

Suppose one line is wrong.

Instead of generating whole program, mask that part:

\[
\text{prefix} \quad <MASK> \quad \text{suffix}
\]

and let model fill the missing part.

CodeBERT can be used for replacing the masked tokens.

So roughly:

\[
\text{buggy line}
\rightarrow
\text{mask}
\rightarrow
\text{model}
\rightarrow
\text{replacement}
\]

This is useful because we already know roughly where
the bug is.

\vspace{8pt}

There are different ways of asking the model to
generate the repair.

\begin{itemize}[leftmargin=18pt,itemsep=2pt]
    \item complete function generation
    \item correct code infilling
    \item single line generation
\end{itemize}

Complete function = regenerate the whole function.

Infilling = keep prefix and suffix, replace the
buggy part.

Single line = only make the required one-line change.

\vspace{8pt}

For infilling:

\[
\boxed{
prefix + [buggy\ location] + suffix
}
\]

The suffix is useful too.

It gives the model information about what comes
after the buggy statement, so it has more context
while predicting the replacement.

\vspace{8pt}

\textbf{many patches}

One generated patch is not enough.

Generate multiple possible patches for the same bug.

\[
P_1,P_2,P_3,\ldots,P_n
\]

then validate them.

Only patches which pass the required checks are
potential repairs.

\vspace{7pt}

Sampling controls how different the generated patches
can be.

Temperature affects the randomness of sampling.

higher temperature $\rightarrow$ more variation

lower temperature $\rightarrow$ more predictable output

Need some balance because too much randomness can
give useless patches.

\vspace{7pt}

Another thing recorded = entropy of generated patch.

Basically based on the probabilities assigned by
the model.

For a token with probability $p$,

\[
H=-\log p
\]

For a complete generated sequence:

mean entropy = average over generated tokens.

sum entropy = total over the sequence.

Lower entropy means the model considers the generated
sequence more natural / likely.

So patches with lower entropy can be tried first.

\[
\text{lower entropy}
\rightarrow
\text{higher priority}
\]

\vspace{8pt}

\textbf{why validation is needed}

LLM can generate code which looks correct but is
actually wrong.

So generated patch $\neq$ accepted patch.

Need to compile / test / verify the patch.

\[
\text{LLM}
\rightarrow
\text{candidate patches}
\rightarrow
\text{validation}
\]

Only then can we say a patch works for the checked
cases.

\vspace{8pt}

\textbf{model size}

Larger models generally gave a higher ratio of
correct patches to generated candidates.

But model size also has a cost.

\[
\text{model size}\uparrow
\Rightarrow
\text{generation time}\uparrow
\]

So there is a tradeoff between quality and speed.

\vspace{7pt}

Also interesting:

models using correct-code infilling and single-line
generation produced a higher ratio of correct fixes
than complete-function generation.

Reason makes sense because the model has a much
smaller thing to change.

\[
\text{small targeted change}
\quad vs \quad
\text{rewrite whole function}
\]

Less unnecessary code to generate.

\vspace{8pt}

The model trained especially on code can perform
well for this task.

So parameter count alone is not everything.

Training data / what the model was trained for matters.

\vspace{8pt}

\textbf{errors}

Two types to think about:

syntax errors

\[
\text{generated code doesn't parse / compile}
\]

semantic errors

\[
\text{code runs but behaviour is wrong}
\]

Correct code infilling had relatively low syntactic
errors and also low semantic errors in the reported
results.

\vspace{8pt}

Traditional APR tools can also be combined with LLMs.

For example:

\[
\text{APR}
+
\text{LLM generated patches}
\]

APR can provide candidate patch ingredients /
locations, while the LLM can generate different
possible repairs.

Then verification decides which candidates work.

\vspace{8pt}

So overall:

\[
\text{buggy program}
\rightarrow
\text{bug location}
\rightarrow
\text{LLM generates patches}
\]

\[
\rightarrow
\text{rank / select candidates}
\rightarrow
\text{validate}
\]

Main idea = don't depend on one generated answer.

Generate several possibilities, use the model's
probabilities to prioritize them, and use program
validation to check the actual repair.

\end{document}